% !TeX program = pdflatex
% ALIUS Bulletin native reconstruction source.
% Compile from the ALIUS-bulletin project root. This file does not include a pre-existing article PDF.
\providecommand{\ALIUSRootPrefix}{}

\ifdefined\ALIUSIssueBuild
\else
\documentclass[a4paper,11pt]{article}
\usepackage[margin=22mm]{geometry}
\usepackage[T1]{fontenc}
\usepackage[utf8]{inputenc}
\usepackage[hidelinks]{hyperref}
\usepackage[protrusion=true,expansion=false]{microtype}
\usepackage{parskip}
\fi
\title{The anthropology of mind: exploring unusual sensations and spiritual experiences across cultures}
\author{Tanya Luhrmann Interviewed by Martin Fortier}
\date{ALIUS Bulletin 1}

\ifdefined\ALIUSIssueBuild
\else
\begin{document}
\fi
\maketitle
\section*{Metadata}
\textbf{Piece type:} interview\par
\textbf{DOI:} 10.34700/7ppf-h459\par
\textbf{Keywords:} voice-hearing, kataphatic prayer, absorption, cognitive penetrability, anthropology of mind\par
\section*{Abstract}
This interview with Tanya Luhrmann examines the cultural and psychological mechanisms through which the supernatural becomes experientially real, with attention to voices, visions, and unusual sensory experiences. Luhrmann describes her trajectory toward an anthropology of mind concerned with how reality acquires its texture for individuals and how that texture shifts at the edges of experience. The conversation turns to participant observation and the role of first-hand anomalous experience in fieldwork on altered states, then to how imagination becomes perception-like among evangelicals through kataphatic prayer. Luhrmann weighs three candidate mechanisms: metacognitive retagging of inner events along the lines of Richard Bentall and Marcia Johnson, training-induced increases in sensory content and imagery vividness, and attentional cultivation of the inner senses, arguing that they can work in concert through micro-moments of attention. She treats the internal-external divide as a continuum rather than a binary and distinguishes the shift from imagination to perception from the sense of reality itself, illustrating the latter through depersonalization. Further sections address the absorption scale, its relation to hypnotizability, dissociation, and psychosis, its genetic component, and the role of personal proclivity in religious practice. The conversation engages the cognitive penetrability debate through a three-level model of theories of mind, habits of mind, and perception, and presents her cross-cultural project across Ghana, China, Thailand, Vanuatu, and the United States on how local theories of mind shape supernatural experience. It closes on participation and Lucien L\'evy-Bruhl's mystical mode of thought, and on new scale items capturing sensorial engagement with a responsive world.\par
\section*{Extracted Text}
\begingroup
\small
\noindent 25 ALIUS Bulletin n1 (2017) aliusresearch.org/bulletin The anthropology of mind Exploring unusual sensations and spiritual experiences across cultures An interview with Tanya Luhrmann By Martin Fortier Tanya Luhrmann luhrmann@stanford.edu Department of Anthropology Stanford University, USA Martin Fortier martin.fortier@ens.fr Institut Jean Nicod ENS/EHESS, Paris, France Citation: Luhrmann, T., \& Fortier, M. (2017). The anthropology of mind: exploring unusual sensations and spiritual experiences across cultures. An interview with Tanya Luhrmann. ALIUS Bulletin, 1, 25-36. To begin, could you please tell us a few words as to how you position yourself within the academic field? In Persuasions of the Witch's Craft (Luhrmann, 1991) you advocated the psychological anthropology approach (1991, pp. 15 et sq.) while at the time acknowledging the importance of cognitive science (1991, pp. 13-14). Since then, you have constantly been reading and exchanging with scholars championing very different approaches-including psychological anthropology, phenomenological anthropology, cognitive anthropology, philosophy, evolutionary psychology, psychiatry, religious studies, developmental psychology, cultural psychology, neuroscience, etc. Besides, you are currently the Principal Investigator of a large interdisciplinary and cross- cultural research project where anthropologists are intensely collaborating with psychologists. How would you define your work today? What are the lines of research you find particularly inspiring? What are your likes and dislikes? I define myself as an anthropologist with deep interests in psychology, and a commitment to using psychological methods to explore the questions ethnography cannot answer. The broadest question that interests me is how the world becomes real for people, particularly when what is real for them seems unreal to others. I'm interested in the texture of reality, and the way it changes for people, and in the quality of individual experience. I want to know about the differences between people, and the way that individual experience shifts and slides. The clearest way to see these moments of difference are in the edges of experience: in voices, visions, the world of the supernatural and the world of psychosis. What drives those experiences has a lot to do with the way people make judgements about what we call mental events. You might call me an anthropologist of mind. While studying magicians in contemporary England (Luhrmann, 1991) or evangelicals in the US (Luhrmann, 2012), it seems that participant observation proved instrumental T. Lurhmann - The Anthropology of Mind ALIUS Bulletin n1 (2017) aliusresearch.org/bulletin 26 for you. Unlike many anthropologists who prefer not to say much about their fieldwork experience, you have written quite openly about the anomalous experiences that you had during your fieldwork (Luhrmann, 1991, p. 319; Luhrmann, 2012, pp. 191-192). According to you, what status should be given to participant observation within anthropology? How native should anthropologists go in order to properly understand their object of investigation? Do you consider first-hand experience to be critical for the study and understanding of altered states of consciousness? I don't think that first-hand experience is essential for understanding human experience-that leads us down the Winch/McIntyre rabbit hole and to the question of whether anything-God, for sure, being American, even owning cattle-can be understood from the outside. And yet there is no question that to have an insider's experience gives you a level of insight that you cannot have any other way. I remember one of my Cambridge supervisors, Stephen Hugh-Jones, telling me he'd gone off to do fieldwork in Amazonia at a time when there was much discussion about why people believed in their gods, and about the more abstract topic of social representation. He took ayahuasca-and he saw spirits. (First, to be clear, he saw London double decker buses.) That changed for him forever the answer to the question of why people believed in spirits. Even if you do not believe that what you see under the influence of ayahuasca is real, it changes your understanding of the event to have some grasp of its phenomenological quality. It adds subtlety to your explanation. In my own case, I had gone into the field anticipating that my explanation for why apparently rational people believed apparently irrational beliefs would be restricted to the realm of narrative and interpretation. When I had an anomalous experience myself-and when I realized that others had these experiences as well-I realized that talk of cognitive interpretation was not enough to capture what I had experienced. That changed the way I thought about religion and what a scholar should strive to explain. In your work, you are especially interested in understanding the mechanisms through which imagination gradually becomes perception-like. In the case of evangelicals, for instance, you describe how they typically start by imagining God, and how, as they train themselves and pretend to be actually interacting with God, their experience becomes more and more real up to the point that they eventually hear God talking directly to I don't think that first-hand experience is essential for understanding human experience... yet there is no question that to have an insider's experience gives you a level of insight that you cannot have any other way. " " T. Lurhmann - The Anthropology of Mind ALIUS Bulletin n1 (2017) aliusresearch.org/bulletin 27 them. Understanding how God comes to be experienced as real is a very intricate question. In your research work, you explore several potential explanations: (1) First, the shift from imagination to perception could be explained in metacognitive terms. For example, Richard Bentall has shown that schizophrenic hallucinations might be explained not by a change in the sensory content itself, but rather by a change in how the content is "tagged" by the mind: the same (more or less sensory-loaded) content can be tagged as "externally generated" - in which case it is experienced to be a real percept-or as "internally generated," in which case it is experienced to be just imagination) (Bentall, 1990; Bentall, Baker, and Havers, 1991). Drawing upon this line of research, you suggest that the shift from imagination to perception could be mediated by a change in metacognition (Luhrmann, 2011a, pp. 72-73; Luhrmann, Nusbaum, and Thisted, 2013, pp. 171-172; Luhrmann et al., 2015b, p. 658). (2) Second, you point out that the training evangelicals perform seems to affect the very content of their sensory experience, and not only how this content is being tagged. That is, the difference between the fictional God that is imagined and the real God that is encountered would be not just metacognitive but would also involve a properly sensory dimension. As you have explained, prayer practice "increases imagery vividness" and "lead[s] to reports of unusual sensory experiences and to reports of unusual sensory experiences associated with the religious ideas" (Luhrmann, Nusbaum, and Thisted, 2013, p. 172). (3) Another proposal is that attention plays an important role in the shift from imagination to perception: "imagery rich practices may make what is imagined more real, not simply because increased attention leads to increased salience, but because the increased attention leads subjects to experience images as more "real"-more percept-like" (Luhrmann, Nusbaum, and Thisted, 2013, p. 161). Now, it could be argued that these three strands of explanation are not as complementary as it may seem at first blush. If the shift from imagination to perception is explained by metacognitive processes (explanation (1)), this means that the experienced content is the same before and after spiritual training and what significantly changes is only how the content is tagged, as opposed to how the content becomes increasingly sensory-loaded. So, arguably, (1) and (2) are not intertwined. As regards explanation (3), attention could here be broadly interpreted as a kind of top- down process modulating access consciousness (Dehaene et al., 2006). One may then reason that through attentional training, a previously unconscious content could suddenly become conscious, which would explain why a rich sensory content is being reported (see (2)). Alternatively, the increase of sensory content could be construed as an instance of decreased sensory gating; as illustrated by the case of schizophrenic T. Lurhmann - The Anthropology of Mind ALIUS Bulletin n1 (2017) aliusresearch.org/bulletin 28 patients, such a decrease typically leads to the experience of sensory overload (Micoulaud-Franchi et al., 2014). While (1) does not seem to fit so well with (2) or (3), on the other hand, it seems reasonable to surmise that (2) and (3) are possibly working complementarily. What is your view on this debate? What do you think is instrumental in the shift from imagination to perception? Do you think that only one of the three aforementioned factors ((1), (2) and (3)) contributes to the shift or that all of them are required in order for the shift to occur? Finally, do you think that the shift from imagination to perception perfectly overlaps the shift from the sense of unreality to the sense of reality? Or do you think that these two pairs of concepts can be orthogonal, and that there is more to the sense of reality than simply experiencing something as being perception-like? These are deep and complicated questions. I tend to assume that the metacognitive tagging changes the sensory qualia of the event post facto, through a micro-moment of attention. That is, the micro-moment decision to infer that the event (some string of words in the mind) is the memory of an event that took place in the world, rather than an event generated by the mind, shifts the experience of the event into a more sensory register. One remembers the event as more external and more sensory. That is more consistent with Richard Bentall's interpretation than Marcia Johnson's, but there is some suggestion in both sets of data that supports this interpretation. Meanwhile, Johnson's original work suggested that the more sensory content in the event under consideration, the more likely the event is to be interpreted as having a source in the external world and thus, more likely to have a sensory quality. So I see (3) as ultimate cause; and I do think that (1) and (2) work in concert. The increased sensory attention of prayer and absorption may lead people to infuse their events with more sensory information, and that in turn may lead to a greater likeliness of a judgment that the event had an external source and thus an experience with a richer sensory trace. I don't think that all three are necessary, but they can work together. I also do not assume that these are the only processes at work, and the only processes in play for any hallucination-like event. At this point, the theoretical model I turn to explains how such phenomena might emerge out of ordinary cognitive process. Imagery and perception depend on many of the same neural structures, as Kosslyn among many others has shown. Increased attention to mental imagery should thus have some effects on a range of image-related cognitive processes: on perceptual processing, on the use of imagery, on unusual sensory experience, and on the vividness of imagery itself-as, indeed, my research has found. The individual trait of absorption (which seems to predispose people to having these unusual experiences may be capturing a similar attention to mental imagery, as many items seem to involve an interest in inner T. Lurhmann - The Anthropology of Mind ALIUS Bulletin n1 (2017) aliusresearch.org/bulletin 29 imagery. Absorption is robustly and significantly correlated with the subjective experience of mental imagery vividness. The puzzle in here, for me, is dissociation/hypnosis. There is a complex and poorly understood relationship between mental imagery vividness, absorption, hypnosis and dissociation. The absorption scale was developed as a pen-and-paper measure of hypnotizability, and while it correlates only modestly (if significantly) with the Stanford Hypnotic Susceptibility Scale, absorption is clearly related to hypnosis. Hypnosis practice increases imagery vividness, and intense spiritual practices can often be described as dissociation-inducing. There is already an active debate in which scholars have argued that most or all voice-hearing experiences are fundamentally related to dissociation due to past trauma. It may be that the pattern and pathway of voice-hearing for those with psychosis differs for those who dissociate and those who do not-regardless of a history of trauma. That should make us rethink some of our assumptions about psychotic hallucination. I increasingly assume that the distinction between internal and external is a continuum rather than a binary. That is certainly what I hear from talking to people with psychosis and those who do not have psychosis but who have unusual experiences: people may say, I know it was not in my head, but I am not sure whether I heard it with my ears. I think certain kinds of events move up and down on the continuum with more ease than others. Indeed, that is the puzzle of the reality monitoring story (Marcia Johnson, Richard Bentall, Yoram Bilu and others)-that there are so few anomalies, as it were, in the way people judge the origin of events. I do not think we have a good account of the infrequency of hallucination-like events. My own sense is that we are constantly having somewhat chaotic events that we correct without being conscious of the corrections. I do think that the sense of what is real is not merely to do with perception, but is another kind of process-related, of course, but not identical. One can see that simply through the fact that perception may be intact/normal while the sense of reality may be deeply disturbed, as in depersonalization or derealization. The sense of reality seems less a judgment about perception, and more a relationship to the act of perception. We look to Martin Fortier to explain our sense of reality in the years to come. In your research, you have demonstrated that the shift from imagination to perception varies widely between individuals. More specifically, personal proclivity for absorption seems to be the driving force of an individual's propensity to have unusual experiences and to successfully transform imagination into percepts (Luhrmann, Nusbaum, and Thisted, 2010). This finding implies that religious experiences are underlain by training (for example, how much kataphatic praying one has performed) but that the T. Lurhmann - The Anthropology of Mind ALIUS Bulletin n1 (2017) aliusresearch.org/bulletin 30 effectiveness of this training is modulated by idiosyncratic characteristics (i.e., one's score on the absorption scale). According to you, what role has this dichotomy between cultural training (praying, meditation, fasting, etc.) and personal disposition (score of the absorption scale) played in the development and shaping of religion across cultures? For example, could it be that people rating high on the absorption scale were preferably selected in imagistic religions, as opposed to doctrinal ones (Whitehouse, 2000)? Do you think that, in some contexts, rituals or trainings are intense and strong enough to make the distinction between highly absorptive individuals and lowly absorptive ones less important (think, for instance, of shamanistic rituals involving the intake of a powerfully hallucinogenic substance)? It has recently been posited that proclivity for absorption is largely underpinned by genetics (Ott et al., 2005). In your opinion, how important is this finding for scholars of religion? Again, an excellent-and substantial-question. I used to think that membership in religious practice was unrelated to personal proclivity. After all, most religious practices have so many members (perhaps a quarter of all Americans are charismatic Christians) and people go to church for such varied reasons (proximity, spouse's preference, etc) that I assumed that differences in proclivity would wash out. I increasingly think proclivity plays some role in the choice of religion. I would expect that role to be more pronounced in smaller religions that require more effort to join. Is proclivity inherently limiting? That is, should someone with low absorption give up an ambition to know God in a sensorially rich manner? Well, no. In my ethnographic work, I have seen people with low absorption, or perhaps better to say a low-absorption orientation to their world, develop the capacity to experience God vividly in a sensorial manner. I think that it is possible to train that style of attentiveness and engagement. At the same time: when I have seen this, I have also wondered about different kinds of proclivity which seem more akin to psychosis. The relationship between ordinary spiritual hallucination-like events and psychosis is fraught and contested (see the just published issue of Schizophrenia Bulletin). For that matter, the question of what processes are involved in psychosis is deeply No great religion has been founded by someone without voices and visions of some sort, even though the predisposition to accept the plausibility of invisible others may be far more broadly distributed, part of the orientation of our evolved brains. " " T. Lurhmann - The Anthropology of Mind ALIUS Bulletin n1 (2017) aliusresearch.org/bulletin 31 contested. We feel increasingly confident that psychosis is complex. I think there may be different processes involved, and that some of those processes may also be involved in spiritual responsivity. Should it matter to scholars of religion that absorption has a genetic component? No great religion has been founded by someone without voices and visions of some sort, even though the predisposition to accept the plausibility of invisible others may be far more broadly distributed, part of the orientation of our evolved brains. And no one, I think, has ever assumed that the people with visions and voices powerful enough to persuade others of their truth were run of the mill individuals. And yet: to begin to figure out what sets highly religious people apart is deeply interesting and, I think, important. The question of cognitive penetrability (i.e., the question of knowing whether high-level beliefs can affect low-level perceptual processes) has sparked much discussions (e.g., Zeimbekis and Raftopoulos, 2015). In your own work-both past and present (Luhrmann, 2012; Luhrmann, 2011a; Luhrmann, 2011b; Cassaniti and Luhrmann, 2014; Luhrmann et al., 2015a; Luhrmann et al., 2015b)-you argue that local models of the mind affect how people experience things. It is not clear, however, whether in doing so you endorse a strong cognitive penetrability thesis, since when you speak of local models of the mind you are referring not only to reflective theories about the mind but also to more procedural processes such as attention. To clarify things a bit, three levels could be distinguished: (1) explicit theories of the mind (e.g., claiming that the mind is made of three components: reason, spirit and soul); (2) habits of the mind (e.g., being good at controlling one's own attention as opposed to being constantly mind-wandering); (3) experience and perception (e.g., hearing someone talking or feeling pain in the back). To illustrate how these levels interact with each other, we could say, for example, that in Buddhism, theories about the mind (the philosophical theory of the five aggregates) does not directly affect perception, whereas what does affect perception is the daily training of the mind (what I have just defined as level (2)). And yet, the training of the mind (i.e., level (2)) seems to be improved and even enhanced by possessing and understanding the theory of the five aggregates (i.e., level (1)). So, in this specific case, it could be tempting to say that level (1) influences level (2) and that level (2) influences level (3), but that (1) does not directly influence (3). What is your view on cognitive penetrability? What does your ethnographic and experimental work suggest about the interplay between explicit theories of the mind, habits of the mind and experience? T. Lurhmann - The Anthropology of Mind ALIUS Bulletin n1 (2017) aliusresearch.org/bulletin 32 My own orientation is to presume that there is indeed downward influence, both from theological orientation and from local cultural expectations, and perhaps upward influence from habits of mind to explicit theories. We now have a large project focused on exploring exactly this topic. That project sets out to understand how cultural variation in ideas about the mind shapes the way people seek and experience the supernatural. We hypothesize that different cultural understandings of the mind-specifically, how separate the mind is from the world, how important inner experience is held to be, and how real the imagination is held to be-shape the way people pay attention to and interpret events they deem supernatural. We propose that although belief in supernatural agents may build upon psychological biases in human cognition, faith is culturally constituted. We are working in five different countries: Ghana, China, Thailand, Vanuatu/Oceania and the US, examining four populations per country: urban charismatic Christian; rural charismatic Christian; urban non-Christian; rural non-Christian. We compare these four populations not only to have comparable groups but also to investigate the impact of charismatic Christianity and industrialization on the way people think about thinking and their experience of the supernatural. In my previous work, I have seen that persons new to a church that taught them that God spoke inside the mind learned to experience thoughts which they might once have treated as self-generated as other-generated (God-generated). I and my colleagues found that prayer practice associated with inner sense cultivation (deliberate attention to inner experience) led those praying to experience what they called God as more person-like and more present, to feel that their inner sensory world became more vivid, and to increase the likelihood that they would experience what they identified as God's voice, visions, and other unusual sensations. The ethnographic literature makes it clear that many religious practices involve training the mind (in particular, cultivating the inner senses) and that training alters the mental experiences of the person trained. The active discussions about embodiment in the psychological literature suggest that these ethnographic findings are an expression of the way that different practices shape subjective experience. Anthropologists have found that different faith interpretations about the believer's mental state have significant consequences for the believer. In particular, there is some evidence that Christianity may change the way that converts understand thoughts and thinking. In an ethnographic analysis of recent converts to Anthropologists have found that different faith interpretations about the believer's mental state have significant consequences for the believer. " " T. Lurhmann - The Anthropology of Mind ALIUS Bulletin n1 (2017) aliusresearch.org/bulletin 33 Pentecostalism in Melanesia, Joel Robbins detailed a sharply increased sense of the social importance of internal states by converts who were unable to rid their minds of thoughts they felt to be sinful. Webb Keane argues in a study of Indonesian Christianity that the expectation of sincerity and inner purity significantly shifted Sumbanese experience of thought and language. Ben Purzycki and Rich Sosis found that different understandings of God's mind also have consequences for the way that people think about human minds. Meanwhile, philosophers, psychologists, and anthropologists have long argued that education and industrialization affects the way people think about minds and mental process (as in, for instance, the work of Karl Popper, Michael Cole and Sylvia Scribner, Ernest Gellner, Robin Horton, and Charles Taylor). Our plan is to conduct both ethnographic and psychological work in order to tease out the relationship between explicit theology, cultural habits of mind, and perception. I do agree that level (1) does not influence level (3) directly. One of your present interests deals with participation, sensuous being-in-the-world, and the feeling of presence. This constellation of themes is evocatively illustrated by, among others, David Abram's (1997) book on the phenomenological approach of sensuous presence and Eduardo Kohn's (2013) book on iconic perception of the world among Amazonian indigenous people. How does this line of research relate to those that you have been exploring so far? Do you think researchers in religious studies and in the cognitive science of religion could gain valuable insights from the rehabilitation of Levy-Bruhl's concept of participation? I do think that researchers could learn something from Levy-Bruhl: I think he is describing a state of being in the world which is far more common, and far more important to religion, than we often realize. Levy-Bruhl saw that what it is to be religious is to experience the world as responsive and full of meaning. In How Natives Think, Levy-Bruhl argued that the distinctive feature of the "primitive" mind was that such primitives experienced themselves as participating in the external world, and the external world as participating in their minds and bodies. Levy-Bruhl called such an orientation "mystical" and he described it as governed by "the law of participation" in which objects are "both themselves and other than themselves." At the end of his life, in the posthumous Notebooks, Levy-Bruhl abandoned the claim that so-called primitive minds were fundamentally different from those of Europeans. He abandoned the term 'prelogical' (1975 [1949]: 99) and began to write of participation as common to all people, different modes of thought rather than different minds. The mystical mode of thought was both affective and conceptual, T. Lurhmann - The Anthropology of Mind ALIUS Bulletin n1 (2017) aliusresearch.org/bulletin 34 and had those features which he had attributed to participation all along: independence from ordinary space and time, logical contradictions (an object is both here and there), identity between objects and their arbitrary features (like hair cuttings and the person from whom they came), "the feeling of a contact, most often unforeseen, with a reality other than the reality given in the surrounding milieu." He thought that the mystical mode intermixed with everyday thought continually in our minds. For him, the puzzle became, "How does it happen that these "mental habits" make themselves felt in certain circumstances and not in others?" I believe that what the absorption scale captures is an interest in feeling sensorially engaged with a responsive world, and that this interest facilitates a sense that the world is alive, aware, intelligent, interested-that it contains, in short, an invisible other. In recent years, I have been developing an alternate scale to the Tellegen absorption scale (the original scale is under copyright protection, which makes it irritatingly difficult to use as a research instrument). Here are some new items that strongly correlate with the absorption scale: "Sometimes the world seems intensely present to me" "When I walk through a forest, I like to think that the trees are murmuring words of wisdom for me" "I have the distinct sense of a wise watchful presence" "When I hear the wave lap against the shore, I sometimes think of how much those waves might know" T. Lurhmann - The Anthropology of Mind ALIUS Bulletin n1 (2017) aliusresearch.org/bulletin 35 References Abram, D. (1997). The Spell of the Sensuous: Perception and Language in a More-Than-Human World. New York: Vintage Books. Bentall, R. (1990). The illusion of reality. A review and integration of psychological research on hallucinations. Psychological Bulletin, 107(1), 82-95. Bentall, R., Baker, G., \& Havers, S. (1991). Reality monitoring and psychotic hallucinations. The British Journal of Clinical Psychology / the British Psychological Society, 30(3), 213- 222. Cassaniti, J., \& Luhrmann, T. (2014). The cultural kindling of spiritual experiences. Current Anthropology, 55(S10), S333-S343. Dehaene, S., Changeux, J.-P., Naccache, L., Sackur, J., \& Sergent, C. (2006). Conscious, preconscious, and subliminal processing: a testable taxonomy. Trends in Cognitive Sciences, 10(5), 204-211. Kohn, E. (2013). How forests think: Toward an anthropology beyond the human. Berkeley/Los Angeles: University of California Press. Luhrmann, T. (1991). Persuasions of the Witch's Craft: Ritual Magic in Contemporary England. Cambridge MA: Harvard University Press. Luhrmann, T. (2011a). Hallucinations and sensory overrides. Annual Review of Anthropology, 40(1), 71-85. Luhrmann, T. (2011b). Toward an anthropological theory of mind. Suomen Antropologi: Journal of the Finnish Anthropological Society, 36(4), 5-69. Luhrmann, T. (2012). When God talks back: Understanding the American evangelical relationship with God. New York: Alfred Knopf. Luhrmann, T., Nusbaum, H., \& Thisted, R. (2010). The Absorption Hypothesis: Learning to Hear God in Evangelical Christianity. American Anthropologist, 112(1), 66-78. Luhrmann, T., Nusbaum, H., \& Thisted, R. (2013). Lord, Teach Us to Pray : Prayer Practice Affects Cognitive Processing. Journal of Cognition and Culture, 13(1-2), 159-177. Luhrmann, T., Padmavati, R., Tharoor, H., \& Osei, A. (2015a). Differences in voice-hearing experiences of people with psychosis in the USA, India and Ghana: interview-based study. The British Journal of Psychiatry, 206(1), 41-44. Luhrmann, T., Padmavati, R., Tharoor, H., \& Osei, A. (2015b). Hearing Voices in Different Cultures: A Social Kindling Hypothesis. Topics in Cognitive Science, 7(4), 646-663. Micoulaud-Franchi, J.-A., Hetrick, W., Aramaki, M., Bolbecker, A., Boyer, L., Ystad, S., ... Vion-Dury, J. (2014). Do schizophrenia patients with low P50-suppression report more perceptual anomalies with the sensory gating inventory? Schizophrenia Research, 157(1-3), 157-162. T. Lurhmann - The Anthropology of Mind ALIUS Bulletin n1 (2017) aliusresearch.org/bulletin 36 Ott, U., Reuter, M., Hennig, J., \& Vaitl, D. (2005). Evidence for a common biological basis of the absorption trait, hallucinogen effects, and positive symptoms: Epistasis between 5HT2a and COMT polymorphisms. American Journal of Medical Genetics Part B (Neuropsychiatric Genetics), 137B(1), 29-32. Whitehouse, H. (2000). Arguments and Icons : Divergent Modes of Religiosity. Oxford: Oxford University Press. Zeimbekis, J., \& Raftopoulos, A. (Eds.). (2015). The cognitive penetrability of perception: New philosophical perspectives. London/New York: Oxford University Press.\par
\endgroup

\ifdefined\ALIUSIssueBuild
\clearpage
\expandafter\endinput
\fi
\end{document}
